% PAGES: 279-279
% Chunk c2 (assigned 279-290) begins here. ch09/ was empty when I started (c1,
% assigned 267-278, had not yet delivered any files), so I could not read a
% c1 file directly; I instead rendered PDF page 278 (folio 262, c1's last
% assigned page) myself via CropBox at 300 dpi to check the boundary. Page 278
% ends with the three-line derivation of $w_{min}$ closing at "$=
% (1-w_{min,cash})w_T,$" tagged (9.39) -- a numbered multi-line display that is
% fully closed on that page, not carried over. Page 279 (this file) opens a
% fresh sentence, "since, from (9.34), it follows that", with its own
% self-contained unnumbered display. Nothing is left open across the c1/c2
% boundary; no \begin{...} needs to be closed at the top of this file.
%
% Notation check (folio 263, 600+ dpi): subscripts on $\Sigma$ and $\mu$ here
% are plain (unbolded) -- $\Sigma_R$, $\bar\mu$ -- since these pages never
% subscript by the bold-lowercase asset-vector $\bfx$; the vector of ones is
% consistently bold ($\bfone$), matching the preamble's note that \bfone is
% the only symbol this book actually bolds in Chapter 9.

since, from (9.34), it follows that
\[
(\mu_P-r_f)\frac{\bfone^t\Sigma_R^{-1}\bar\mu}{\bar\mu^t\Sigma_R^{-1}\bar\mu} \;=\;
1-w_{min,cash}.
\]

From (9.38) and (9.39), we conclude that the minimum variance portfolio with expected
return $\mu_P$ is obtained for the following asset allocation:
\begin{itemize}
\item weight $w_{min,cash}$ of the cash position given by
\begin{equation}
w_{min,cash} \;=\; 1-\frac{\mu_P-r_f}{\bar\mu^tw_T};
\end{equation}
\item asset weights vector $w_{min}$ given by
\begin{equation}
w_{min} \;=\; (1-w_{min,cash})\,w_T.
\end{equation}
\end{itemize}

The pseudocode for the asset allocation for the minimum variance portfolio computed
using the formulas (9.40) and (9.41) can be found in Table 9.2.

We revisit below the example from section 9.3 and find the minimum variance portfolio
by using the tangency portfolio.

\par\medskip\noindent\textit{Example:}\ Find the \$100 million minimum variance
portfolio with expected return equal to 2.25\% and invested in cash and four assets
with expected returns equal to 2\%, 1.75\%, 2.5\%, and 1.5\%, and with the following
covariance matrix of their returns:
\[
\begin{pmatrix}
0.09 & 0.01 & 0.03 & -0.015 \\
0.01 & 0.0625 & -0.02 & -0.01 \\
0.03 & -0.02 & 0.1225 & 0.02 \\
-0.015 & -0.01 & 0.02 & 0.0576
\end{pmatrix}.
\]

Assume that the risk free rate is 1\%.

\begin{answerx}
% Left OPEN at the end of this file: the derivation continues onto PDF page
% 280 (folio 264), which closes it with "$w_{min} = ... $" followed by the
% square. sec09_c2_2.tex must NOT reopen \begin{answerx}.
Let $\mu_1 = 0.02$, $\mu_2 = 0.0175$, $\mu_3 = 0.025$, $\mu_4 = 0.015$, and $r_f =
0.01$. We look for the assets and cash allocation corresponding to the minimum
variance portfolio with expected return $\mu_P = 0.0225$.

Note that $\bar\mu = \mu - r_f\bfone$ and $\Sigma_R$ are given by
\[
\bar\mu \;=\; \begin{pmatrix} 0.01 \\ 0.0075 \\ 0.015 \\ 0.005 \end{pmatrix}; \quad
\Sigma_R \;=\; \begin{pmatrix}
0.09 & 0.01 & 0.03 & -0.015 \\
0.01 & 0.0625 & -0.02 & -0.01 \\
0.03 & -0.02 & 0.1225 & 0.02 \\
-0.015 & -0.01 & 0.02 & 0.0576
\end{pmatrix}.
\]

The vector $\Sigma_R^{-1}\bar\mu = \text{linear\_solve\_cholesky}(\Sigma_R,\bar\mu)$ is
computed by using the Cholesky decomposition to solve the linear system corresponding
to the matrix $\Sigma_R$ and to the right hand side $\bar\mu$. From (9.36), we find
that the asset weights vector of the tangency portfolio is
\[
w_T \;=\; \frac{1}{\bfone^t\Sigma_R^{-1}\bar\mu}\Sigma_R^{-1}\bar\mu \;=\;
\begin{pmatrix} 0.1587 \\ 0.3660 \\ 0.2650 \\ 0.2103 \end{pmatrix}.
\]

From (9.40), we find that the weight of the cash position in the minimum variance
portfolio is
\[
w_{min,cash} \;=\; 1-\frac{\mu_P-r_f}{\bar\mu^tw_T} \;=\; -0.3356,
\]
